\section{Design}
\subsection{State ownership and interfaces}
Figure~\ref{fig:architecture} identifies the owners of the three decisions.
The workflow runtime accepts a graph of tasks and named data dependencies.
The scheduler maintains logical task states, unresolved dependency counts,
dispatch choices, and retry generations. It does not select a concrete source
replica for every task input during placement. The data controller maintains
DataID records, available replicas, invalidation, retention, and backup state.
Each worker agent maintains its local view of materialized data and performs
input resolution and movement. The generic TaskVine manager remains the sole
owner of the underlying task-manager API.

\begin{figure}[t]
\centering\includegraphics[width=\columnwidth]{architecture.pdf}
\caption{DataVine ownership and gates. Dashed arrows carry completion and
lifecycle information; data transfer does not route through the logical
scheduler. The manager and controller still perform coordination work.}
\label{fig:architecture}
\end{figure}

An input reference denotes data identity rather than a placement-time pathname
on a specific peer. The worker checks for a usable local copy, asks the
controller for a source when needed, and resolves that identity to a physical
input in the task sandbox. Replica location can change between dispatch and
preparation without changing the logical DAG edge. Conversely, a stale source
is not evidence that the producer's logical output should be accepted from a
different, obsolete attempt. Generations bind attempts and their control
messages to the scheduler's current state.

Separation does not remove all per-file work from the workflow runtime.
Completing a logical attempt still updates pending consumers and coordinates
mark-dead operations with the controller. The design changes ownership of
replica decisions, rather than asserting a zero-cost control path. This
qualification matters at high task rates, where serialized manager send,
status handling, and lifecycle operations can become dominant.

\subsection{Dispatch before local readiness}
After a producer's successful physical attempt is accepted, the scheduler
updates logical dependencies and makes newly eligible consumers dispatchable.
The scheduler may assign a consumer before its destination has prepared the
input. The assigned descriptor waits in a bounded worker queue. For the
fork-based DataVine execution path, the worker first checks whether a suitable
execution-library opportunity exists; absent one, it postpones the agent's
preparation operation. Once preparation starts, the agent resolves generated
inputs and initiates required transfers. Execution remains conditional on
input readiness, library availability, and resource matching.

This order prevents a deep descriptor queue from automatically becoming an
equally deep collection of staged sandboxes. It also changes prefetch distance:
earlier preparation can hide a long transfer, whereas later preparation can
reduce unnecessary staging. The correct empirical question is where this
tradeoff changes sign. We therefore implement an eager preparation control
that bypasses only the preliminary execution-opportunity check. It retains
all input-readiness and actual execution checks. The control is not an unsafe
mode that executes without data.

The design does not provide two independently optimized transfer and compute
windows. Preparation and execution remain related through the library gate;
input transfer may occupy an opportunity that could otherwise serve another
call. Decoupling those windows further is a possible extension, but requires
separate bounds on transfer bytes, sender contention, and prepared-data
occupancy. We avoid attributing such a mechanism to this implementation.

\subsection{Data path and retention}
For each generated input, a worker uses a valid local replica when available.
Otherwise the controller resolves a live source; peer data can be transferred
directly between workers. A controller backup is eligible as fallback only
after its admission has completed. This ordering avoids publishing an
incomplete backup as if it were a usable replica. A peer becoming unavailable
can trigger a new resolution rather than forcing the scheduler to encode a
permanent source into the task descriptor.

The default intermediate policy retains live worker copies and performs
controller backup in the background. Controlled performance comparisons can
select worker-local intermediates, disabling this backup while preserving
requested-output handling. These configurations have different recovery
costs and must not be combined into a single speedup claim. Shared-filesystem
experiments in the wider artifact explore a possible alternative data path;
the production worker does not currently choose it through an adaptive
three-way cost model.

Data lifetime is tied to consumers and requested outputs, not merely to the
producer's return. Once the graph can establish that an intermediate is no
longer needed, lifecycle updates permit reclamation. Dynamic graph extension
makes this more subtle: an open workflow can introduce a future consumer, so
retained data may accumulate until sealing establishes a closed consumer set.
A bounded descriptor queue is therefore not a global bound on intermediate
storage. The controller's retention policy and workflow closure remain part
of the capacity model.

\subsection{Worker execution feedback}
The default dispatch capacity is $Q=8C$. The elastic execution window starts
at $W=2C$ and samples worker pressure every 250\,ms. CPU saturation, runnable
pressure, and memory are separate signals. In the current policy, CPU at or
above 95\% together with runnable pressure above $3C$ for two samples reduces
the window by $C$, with a floor of $2C$. The $3C$ value is a feedback threshold,
not a hard bound on concurrent processes. Memory pressure at 80\% causes a
stronger backoff and drain; declared task memory supplies an additional
admission check.

When calls are waiting, low measured CPU (below 85\%) and memory (below
70\%) allow multiplicative growth; otherwise an eligible growth step adds
$C$, capped by configured library slots. The policy holds near its 90\% CPU
target when the runnable population is at least $C$, and inserts a cooldown
after backoff. Before increasing $W$, it estimates a memory-limited window
from the current running count $R$ and measured memory fraction $m$:
\begin{equation}
W_{\mathrm{next}}\leftarrow\min\left(W_{\mathrm{candidate}},
\max\left(R,\left\lfloor0.8R/m\right\rfloor\right)\right).
\end{equation}
This guard applies when $R>0$ and $m>0$. A memory-fraction increase above
0.10 in one sample defers growth to avoid extrapolating from a partially
faulted-in footprint. These thresholds are implementation choices tested
against fixed windows, not tuned optima. Shrinking admission does not kill
running calls: the population drains before newly permitted calls enter.

The policy aims to respond to current aggregate behavior without requiring
accurate runtime predictions for every task class. It is not a proof of memory
safety or stability. A task can allocate more memory than declared between
samples, and low CPU usage does not imply that more processes will help a
memory-bandwidth-bound workload. Fixed windows at $C$, $2C$, and $4C$ therefore
remain essential controls. We report the best tested fixed choice as an
empirical reference, not as an oracle for every possible setting.

\subsection{Recall, failure, and result visibility}
An assigned task that has not begun execution may be recalled under its
current generation. The runtime distinguishes that state from a running
attempt; recall is not process migration. Retrying under a new generation
must prevent delayed messages from an earlier attempt from advancing the
new attempt's state. This is especially important when worker loss,
re-resolution, and completion notification overlap.

Physical success releases logical descendants before background persistence,
but user-visible requested results wait for durable admission. A completed
worker task and a completed workflow result are thus deliberately different
states. Recovery tests exercise worker loss, lifecycle handling, and replay.
The contract does not give arbitrary user side effects exactly-once semantics.
A retried task that writes to an external database needs application-level
idempotence or a stronger transactional interface.

Finally, controller-local storage determines the failure domain. A backup on
the controller's local temporary filesystem survives a worker's loss but is
not a cross-host durable service if that controller host disappears. We make
this assumption explicit rather than generalizing worker-loss tests into a
claim of controller-host fault tolerance.
